\documentclass[11pt]{article}
\usepackage[T1]{fontenc}
\usepackage{lmodern,amsmath,amssymb,amsthm,mathtools}
\usepackage[margin=1in]{geometry}
\usepackage{microtype,booktabs,array,longtable,enumitem,xcolor,listings,xurl}
\usepackage[colorlinks=true,linkcolor=blue!45!black,urlcolor=blue!45!black,citecolor=blue!45!black]{hyperref}
\usepackage{fancyhdr}
\pagestyle{fancy}\fancyhf{}
\fancyhead[L]{Scalar iteration uniqueness for iPRA}\fancyhead[R]{\thepage}
\setlength{\headheight}{14pt}
\setlength{\emergencystretch}{2em}
\setlist{itemsep=3pt,topsep=5pt}
\lstset{basicstyle=\ttfamily\scriptsize,breaklines=true,columns=fullflexible,keepspaces=true,frame=single,rulecolor=\color{black!20},showstringspaces=false}
\newtheorem{theorem}{Theorem}[section]
\newtheorem{lemma}[theorem]{Lemma}
\newtheorem{proposition}[theorem]{Proposition}
\newtheorem{corollary}[theorem]{Corollary}
\theoremstyle{definition}\newtheorem{definition}[theorem]{Definition}
\newcommand{\iPRA}{\mathsf{iPRA}}
\newcommand{\dotminus}{\mathbin{\dot{-}}}
\newcommand{\file}[1]{{\small\nolinkurl{#1}}}
\hypersetup{pdftitle={Scalar iteration uniqueness for iPRA},pdfauthor={Research report prepared for Vladimir Reshetnikov}}
\title{Scalar iteration uniqueness for iPRA\\[5pt]
\large A complete arithmetic alternative and a small validated-checker refinement}
\author{Research report prepared for Vladimir Reshetnikov}
\date{November 11, 2026}
\begin{document}
\maketitle
\begin{abstract}
The induction principle in the current scalar presentation of full
first-order intuitionistic primitive recursive arithmetic can be replaced
by unconditional uniqueness of a one-state iterator.  From a derivation of
$f(Sz)=h(f(z))$ under a fixed context, the new rule concludes
$f(t)=I(t,f(0);u.h(u))$.  It has no induction hypothesis and no separate
base premise.  We prove equivalence for all contextual first-order
sequents, hence whole-domain bi-interpretability and equiconsistency
with iPRA through the existing verified scalar bridge.

The main new proof closes the scalar equality-reflection gap explicitly
left open in the incoming \texttt{scope\_kernel} report.  A modified
finite-state ramp constructs predecessor using uniqueness.  An explicit
scalar accumulator then establishes
$a+(b\mathbin{\dotminus}a)=b+(a\mathbin{\dotminus}b)$ and a faithful
equality test, before pairing or general primitive recursion is used.

The size result is deliberately separate.  A straightforward complete
checker for the new uniqueness calculus is 3393 normalized Python bytes,
versus 3322 for its matched induction control.  Thus this new arithmetic
basis is \emph{not} a demonstrated improvement under the proof-checker-size
criterion.  A four-byte refinement of the existing, unchanged validated
calculus reduces its complete checker from 870 to 866 bytes.  Neither
result is a proof of absolute minimality or a new overall compact-codec
minimum.  The mathematical equivalence is a written proof, not a new
Lean or Rocq formalization.
\end{abstract}
\newpage
\tableofcontents
\newpage

\section{Result, comparison, and provenance}

The task asks for a simpler complete axiomatization, with equiconsistency
and bi-interpretability rather than just agreement about standard natural
numbers.  This report establishes a new complete arithmetic alternative
and records precisely where its simplicity claim stops.  Removing a local
induction hypothesis simplifies the arithmetic inference interface.
It does not, in the supplied implementation, make the complete checker
smaller.  Under the stricter checker-size objective, the mathematical
part of the requested improvement remains unachieved.

The comparison is pinned to ProveIt commit
\begin{center}\small\texttt{11be5aabc6b07262046e2553419bc6b05271d178}.\end{center}
The inspected project is
\file{Logic/PrimitiveRecursiveArithmetic/ConciseIPRA}.
Its README and all six then-present incoming archives were inspected;
the exact audit and source hashes accompany this report.  The current
repository already contains zero-computation normalization, six-rule
presentations, one-term and template-free equality rules, and several
introduction-mode refinements.  These are not new contributions here
\cite{project,zero-normal,template-free}.

The incoming \texttt{scope\_kernel} report proves scalar uniqueness implies
zero-equation induction by a retraction argument, but leaves predecessor
and faithful numerical equality reflection as an open bootstrap problem
\cite{scope}.  We reproduce its retraction argument with attribution and
complete that bootstrap.  Recursion uniqueness as a general replacement
for induction is classical.  For example, Okada and Scott study uniqueness
of iteration in a substantially richer typed setting \cite{okada-scott}.
The claim here is a completion relative to this repository's restricted
scalar language, not historical priority for recursion uniqueness.

\begin{center}
\begin{tabular}{@{}lrrrr@{}}
\toprule
Complete checker & Bytes & Tokens & AST nodes & Gzip-9\\
\midrule
Matched native equation-progress control & 3322 & 1084 & 1141 & 984\\
New scalar uniqueness calculus & 3393 & 1124 & 1175 & 1011\\
Existing full validated lexical codec & 870 & 383 & 442 & 507\\
Same codec with scope reuse & 866 & 381 & 439 & 504\\
\bottomrule
\end{tabular}
\end{center}
All measurements use CPython 3.12.14, \texttt{ast.unparse}, and one final
newline.  Significant tokens omit comments and layout tokens; AST nodes
count the complete parsed module.  Gzip uses level 9 and timestamp zero,
and is informational.  The two pairs have different certificate formats.
Only the first pair implements the new UR arithmetic.  The 866-byte
program retains the old calculus and is not a UR implementation.  The
repository's separately reported 516-byte profile has a different public
syntax contract; no comparison across those interfaces is presented as a
new minimum.

\section{A complete native axiom and rule inventory}
\label{sec:native-inventory}

For a presentation with no ambient logical axioms, take one numerical sort
and terms
\[
 t::=x\mid0\mid St\mid I(n,a;u.h).
\]
Only $u$ in $h$ is bound.  Formulas are
\[
 A::=(s=t)\mid Q_yA\mid D_{x,y}A.
\]
All substitutions avoid capture; contexts are finite reusable lists of
formulas.  The bound coordinates in $D$ are distinct.  The designated
contradiction is $0=S0$.  No classical logical rule is implicit.

These eight patterns form the entire calculus.  Its complete checker is
\file{code/uniqueness/scalar_ur.py}.

\paragraph{Hypothesis.} $\Gamma\vdash A$ when $A\in\Gamma$.

\paragraph{Two computation axioms.} In every context,
\[
 I(0,a;u.h)=a,\qquad
 I(Sn,a;u.h)=h[I(n,a;u.h)/u].
\]

\paragraph{Crossed transport in one term.} For a numerical template $k(v)$,
\[
 \frac{\Gamma\vdash s=t\qquad\Gamma\vdash k(s)=r}
      {\Gamma\vdash r=k(t)}.                            \tag{Tr}
\]
Sources are unrestricted.  In particular, literal zero is allowed.

\paragraph{Conditional-closure elimination.}
\[
 \frac{\Gamma\vdash Q_yB(y)\qquad\Gamma\vdash B(0)}
      {\Gamma\vdash B(t)}.                             \tag{QE}
\]

\paragraph{Paired elimination.}
\[
 \frac{\Gamma\vdash D_{x,y}B(x,y)\qquad
       \Gamma,Q_yB(b,y)\vdash A}
      {\Gamma\vdash A}.                               \tag{DE}
\]
The eigenparameter $b$ is fresh for the retained context, conclusion,
and fixed parameters of the matrix.

\paragraph{Paired introduction by cases only.}
\[
 \frac{\Gamma,B(t,0)\vdash B(t,Sz)}
      {\Gamma\vdash D_{x,y}B(x,y)}.                    \tag{DI}
\]
Here $z$ is fresh for $\Gamma,t$ and all fixed matrix parameters.
This rule treats atomic and compound matrices alike; there is no
progress option and no literal-zero witness dispatch.

\paragraph{Scalar recursion uniqueness.}
\[
 \frac{\Gamma\vdash f(Sz)=h[f(z)/u]}
      {\Gamma\vdash f(t)=I(t,f(0);u.h)}.               \tag{UR}
\]
The template $h(u)$ has no free $z$.  The variable $z$ is fresh for
$\Gamma,t$ and every passive parameter.  It is a fresh numerical
parameter, not a function variable.  In particular, the rule's premise
has no local arithmetic assumption.

There is no further equality axiom, induction axiom, successor axiom,
structural proof instruction, or logical axiom.  Reflexivity follows by
using $I(0,t;u.0)=t$ as both transport premises with $k(v)=v$.
Symmetry follows by transporting $s=t$ along $s=s$ with that same
template.  Transitivity and term congruence then follow directly.
Full formula transport is recovered by structural induction on the finite
formula syntax, using DI, DE, and QE.  Weakening, renaming, substitution,
and assumption cut are admissible by induction on derivations.  UR is
stable under parameter substitution after its eigenparameter is freshened.

\subsection{The ordinary logical reading}
The readings are
\[
 Q_yB:=B(0)\to\forall y\,B(y),\qquad
 D_{x,y}B:=\exists x\bigl(B(x,0)\to\forall y\,B(x,y)\bigr).
\]
Ordinary intuitionistic logic is reconstructed without any use of
equation progress.  For clarity, a compact description of its macros is
given here.  The exact inherited derivations are included in
\file{reference/inherited_choice_basis.tex}, especially the sections on
cases, equality, binder congruence, and the entire ordinary logical
interface \cite{choice}.

First replace a $Q$ formula, when needed inside a macro, by an equivalent
$D$ formula with an unused first coordinate.  Set
$C(v,a,b)=I(v,a;u.b)$.  On equations define a gate
$G_v(s=t)$ as $C(v,0,s)=C(v,0,t)$ and extend it through $D$.
It is provable at zero and equivalent to the original formula at a
successor.  A formula selector $M_v(A,B)$ is built by selecting the two
equation endpoints with $C$, extending through aligned $D$ binders,
and padding the shallower input by $D_{w,y}G_y$.  Thus its zero instance
is equivalent to $A$, and every successor instance to $B$.
Put $p(z)=I(z,0;u.C(u,1,0))$, so $p(1)=1$ and $p(2)=0$ by computation.
With all displayed extra bound variables fresh, the macros are
\begin{align*}
 \bot_*&:=(0=S0),\\
 A\to_*B&:=D_{w,z}M_z(A,B),\\
 A\land_*B&:=D_{w,z}G_zM_{p(z)}(A,B),\\
 \exists_*x\,A&:=D_{x,y}G_yA,\\
 A\lor_*B&:=\exists_*v\,M_v(A,B),\\
 \forall_*x\,A(x)&:=A(0)\land_*D_{w,z}G_zA(z).
\end{align*}
The rules for these macros are precisely intuitionistic natural deduction:
assumption; introduction and elimination for $\land,\lor,\to,\forall,
\exists$; and explosion from $\bot$.  Quantifier introduction and
existential elimination have their usual fresh-parameter conditions.
For $\lor$-elimination both branches must prove the same conclusion.
No rule of excluded middle or double-negation elimination is present.
The numerical case rule is derived using DI with an unused witness,
then DE and QE.  It gives coverage.  Under $0=S0$, selector transport
gives every equation, and external formula-structural induction gives
explosion.  These observations justify using ordinary intuitionistic
reasoning in the arithmetic proof that follows without relying on the
induction rule being replaced.

